% ---------------------------------------------------------
% PREAMBLE (same template as Companion X–XXII)
% ---------------------------------------------------------
\documentclass[11pt,a4paper]{article}

\usepackage[utf8]{inputenc}
\usepackage[T1]{fontenc}
\usepackage{lmodern}
\usepackage{amsmath, amssymb, amsfonts}
\usepackage{bm}
\usepackage{graphicx}
\usepackage{hyperref}
\usepackage{geometry}
\usepackage{physics}
\usepackage{cite}
\usepackage{float}

\geometry{margin=1in}

% ---------------------------------------------------------
% TITLE
% ---------------------------------------------------------
\title{
\textbf{Companion XXIII}\\[4mm]
\textbf{The Formation of the First Black Holes in the}\\
\textbf{Relaxation-Driven Cyclic Cosmology}\\[2mm]
\textbf{Seed Channels, Early Growth, and High-Redshift Quasars}
}

\author{
Michael Lehmann\\
Independent Researcher\\
\texttt{mi.lehmann@gmx.de}
}

\date{April 2026}

\begin{document}

\maketitle

% ---------------------------------------------------------
% ABSTRACT
% ---------------------------------------------------------
\begin{abstract}
The discovery of supermassive black holes (SMBHs) with masses exceeding 
$10^{8}\,M_{\odot}$ at redshifts $z > 10$ poses a major challenge to the standard 
$\Lambda$CDM cosmology. The rapid assembly of such massive objects requires either 
extremely efficient accretion, massive initial seeds, or non-standard early structure 
formation. In the Relaxation-Driven Cyclic Cosmology (RDCC), however, the suppression of 
small-scale power, the delayed collapse of low-mass halos, and the modified cooling and 
fragmentation physics fundamentally alter the pathways for early black hole formation.

In this Companion we develop the first RDCC framework for the formation and growth of the 
earliest black holes. We analyze three seed formation channels—Population~III remnants, 
direct-collapse black holes (DCBHs), and stellar-dynamical runaway collapse—and show how 
relaxation-modified gravitational dynamics shift their efficiency and characteristic mass 
scales. We derive the RDCC-modified accretion rates, merger histories, and early black hole 
mass functions, and we compute the conditions under which SMBHs with 
$M_{\bullet} \gtrsim 10^{8}\,M_{\odot}$ can form by $z \sim 10$.

Building on the early-galaxy framework of Companion~XXII and the baryonic physics of 
Companion~XX, we show that RDCC naturally enhances the formation of massive black hole seeds 
while suppressing low-mass fragmentation. The resulting SMBH population is consistent with 
JWST observations of high-redshift quasars and provides clear predictions for LISA and future 
deep-field surveys.

This Companion establishes early black hole formation as a key observational probe of the 
relaxation dynamics of RDCC and extends the model into the regime of the first quasars.
\end{abstract}

% ---------------------------------------------------------
% SECTION 1 — MOTIVATION AND OVERVIEW
% ---------------------------------------------------------
\section{Motivation and Overview}

The discovery of luminous quasars at redshifts $z > 10$ with inferred black hole masses 
$M_{\bullet} \gtrsim 10^{8}\,M_{\odot}$ presents one of the most significant challenges to the 
standard $\Lambda$CDM cosmology. The rapid assembly of such massive objects requires either 
extremely efficient accretion, unusually massive initial seeds, or non-standard early 
structure formation. None of these conditions arise naturally in the standard model.

In $\Lambda$CDM, the earliest black hole seeds form through:

\begin{itemize}
    \item Population~III (Pop~III) remnants with masses 
          $M_{\bullet} \sim 10$--$100\,M_{\odot}$,
    \item direct-collapse black holes (DCBHs) with 
          $M_{\bullet} \sim 10^{4}$--$10^{6}\,M_{\odot}$,
    \item stellar-dynamical runaway collapse in dense clusters.
\end{itemize}

However, each of these channels faces severe constraints:

\begin{itemize}
    \item Pop~III remnants require sustained Eddington accretion for $\sim 500$ Myr,
    \item DCBHs require rare conditions (strong Lyman–Werner flux, low metallicity),
    \item stellar-dynamical collapse requires extremely high gas inflow rates.
\end{itemize}

Recent JWST observations exacerbate these tensions by revealing:

\begin{itemize}
    \item quasars with $M_{\bullet} \gtrsim 10^{8}\,M_{\odot}$ at $z > 10$,
    \item compact, massive host galaxies at early times,
    \item suppressed low-mass galaxy populations,
    \item strong clustering of early luminous sources.
\end{itemize}

These findings suggest that early black hole formation may proceed differently from the 
standard scenario.

% ---------------------------------------------------------
\subsection*{1.1 RDCC Perspective}

The Relaxation-Driven Cyclic Cosmology (RDCC) modifies early structure formation through the 
scale-dependent relaxation rate $\Gamma(a)$, which suppresses small-scale power and delays the 
collapse of low-mass halos. As shown in Companion~XVIII–XXII, this leads to:

\begin{itemize}
    \item reduced abundance of minihalos,
    \item delayed cooling and fragmentation,
    \item a more top-heavy Pop~III initial mass function,
    \item enhanced formation of massive gas clouds,
    \item smoother early galaxy assembly.
\end{itemize}

These effects directly influence black hole seed formation:

\begin{itemize}
    \item \textbf{Pop~III remnants become more massive} due to larger fragmentation scales.
    \item \textbf{Direct-collapse conditions become more common} due to suppressed $\mathrm{H}_2$ cooling.
    \item \textbf{Stellar-dynamical collapse is enhanced} by higher gas inflow rates in massive halos.
    \item \textbf{Early accretion is more efficient} due to reduced feedback in low-density environments.
\end{itemize}

Thus, RDCC provides a natural pathway to forming massive black hole seeds without requiring 
fine-tuned initial conditions.

% ---------------------------------------------------------
\subsection*{1.2 Goals of This Companion}

The goals of this Companion are:

\begin{itemize}
    \item to derive the RDCC-modified conditions for Pop~III remnant formation,
    \item to compute the RDCC-enhanced direct-collapse black hole (DCBH) channel,
    \item to analyze stellar-dynamical runaway collapse in RDCC halos,
    \item to model early black hole accretion and merger histories,
    \item and to predict the RDCC black hole mass function at $z > 6$.
\end{itemize}

These results build directly on the early-galaxy framework of Companion~XXII and the baryonic 
physics of Companion~XX.

% ---------------------------------------------------------
\subsection*{1.3 Structure of This Companion}

The structure of this Companion is as follows:

\begin{itemize}
    \item Section~2 derives the RDCC-modified seed formation channels.
    \item Section~3 analyzes early black hole accretion and feedback.
    \item Section~4 develops the RDCC black hole mass function and merger history.
    \item Section~5 compares RDCC predictions with JWST quasars and high-$z$ AGN.
    \item Section~6 summarizes observational signatures and provides concluding remarks.
\end{itemize}

Together, these results establish early black hole formation as a powerful probe of the 
relaxation dynamics of RDCC.

% ---------------------------------------------------------
% SECTION 2 — RDCC-MODIFIED BLACK HOLE SEED FORMATION CHANNELS
% ---------------------------------------------------------
\section{RDCC-Modified Black Hole Seed Formation Channels}

The formation of the first black hole seeds marks the beginning of the cosmic black hole 
population and sets the initial conditions for the emergence of high-redshift quasars. In the 
standard $\Lambda$CDM model, three primary channels contribute to early black hole formation:

\begin{enumerate}
    \item Population~III (Pop~III) stellar remnants,
    \item direct-collapse black holes (DCBHs),
    \item stellar-dynamical runaway collapse in dense clusters.
\end{enumerate}

In RDCC, however, the suppression of small-scale power, the delayed collapse of low-mass 
halos, and the modified cooling and fragmentation physics alter the efficiency, mass scale, 
and prevalence of each channel. This section derives the RDCC-modified conditions for seed 
formation and quantifies the resulting seed mass spectrum.

% ---------------------------------------------------------
\subsection*{2.1 Pop~III Remnant Black Holes}

In $\Lambda$CDM, Pop~III stars form in minihalos with masses 
$M \sim 10^{5}$--$10^{6}\,M_{\odot}$ and produce black hole remnants with masses 
$M_{\bullet} \sim 10$--$100\,M_{\odot}$. In RDCC, the fragmentation scale is increased due to:

\begin{itemize}
    \item reduced gas densities,
    \item delayed cooling,
    \item larger Jeans masses,
    \item suppressed small-scale structure.
\end{itemize}

As shown in Companion~XXII, the characteristic Pop~III stellar mass becomes
\begin{equation}
    M_{\star,\mathrm{III,RDCC}}
    \simeq M_{\star,\mathrm{III}}
    \left[
        1 + \chi_{\mathrm{frag}}
        \left(\frac{M}{M_{\mathrm{rel}}}\right)^{-\alpha}
    \right],
\end{equation}
with $\chi_{\mathrm{frag}} \sim 0.1$.

Thus, Pop~III remnants become more massive:
\begin{equation}
    M_{\bullet,\mathrm{III,RDCC}}
    \sim 30\text{--}300\,M_{\odot}.
\end{equation}

Consequences:

\begin{itemize}
    \item fewer low-mass seeds,
    \item more massive initial remnants,
    \item enhanced probability of early Eddington growth.
\end{itemize}

% ---------------------------------------------------------
\subsection*{2.2 Direct-Collapse Black Holes (DCBHs)}

DCBHs form when gas collapses nearly isothermally at $T \sim 8000$ K without fragmenting. 
In $\Lambda$CDM, this requires:

\begin{itemize}
    \item strong Lyman–Werner (LW) flux to suppress $\mathrm{H}_2$,
    \item low metallicity,
    \item high inflow rates ($\dot{M} \gtrsim 0.1\,M_{\odot}\,\mathrm{yr}^{-1}$).
\end{itemize}

In RDCC, these conditions become significantly more common due to:

\begin{itemize}
    \item suppressed $\mathrm{H}_2$ formation (Companion~XXII),
    \item delayed Pop~III enrichment,
    \item larger gas reservoirs in atomic-cooling halos,
    \item reduced fragmentation.
\end{itemize}

The critical LW flux becomes
\begin{equation}
    J_{\mathrm{crit,RDCC}}
    = J_{\mathrm{crit}}
      \left[
        1 - \chi_{\mathrm{LW}}
        \left(\frac{M}{M_{\mathrm{rel}}}\right)^{-\alpha}
      \right],
\end{equation}
with $\chi_{\mathrm{LW}} \sim 0.2$.

Thus, DCBH formation is enhanced by a factor of $\sim 2$–$3$.

The characteristic DCBH seed mass becomes
\begin{equation}
    M_{\bullet,\mathrm{DCBH,RDCC}}
    \sim 10^{5}\text{--}10^{6}\,M_{\odot}.
\end{equation}

% ---------------------------------------------------------
\subsection*{2.3 Stellar-Dynamical Runaway Collapse}

Runaway collapse occurs in dense stellar clusters when the core-collapse timescale is shorter 
than the massive-star lifetime. In $\Lambda$CDM, this requires:

\begin{itemize}
    \item high gas inflow rates,
    \item compact star-forming regions,
    \item low metallicity.
\end{itemize}

In RDCC:

\begin{itemize}
    \item delayed fragmentation increases cluster masses,
    \item reduced cooling leads to more massive gas cores,
    \item smoother early galaxy assembly enhances inflow rates.
\end{itemize}

The characteristic seed mass becomes
\begin{equation}
    M_{\bullet,\mathrm{run,RDCC}}
    \sim 10^{3}\text{--}10^{4}\,M_{\odot}.
\end{equation}

% ---------------------------------------------------------
\subsection*{2.4 Combined RDCC Seed Mass Function}

The RDCC seed mass function is a weighted sum of the three channels:
\begin{equation}
    \Phi_{\bullet,\mathrm{seed}}(M)
    = \Phi_{\mathrm{III}}(M)
    + \Phi_{\mathrm{DCBH}}(M)
    + \Phi_{\mathrm{run}}(M).
\end{equation}

RDCC predicts:

\begin{itemize}
    \item a suppressed low-mass tail ($M_{\bullet} < 30\,M_{\odot}$),
    \item an enhanced intermediate-mass population ($10^{3}$–$10^{4}\,M_{\odot}$),
    \item a more common DCBH channel,
    \item a higher median seed mass.
\end{itemize}

% ---------------------------------------------------------
\subsection*{2.5 Summary}

RDCC modifies black hole seed formation through:

\begin{itemize}
    \item more massive Pop~III remnants,
    \item enhanced direct-collapse conditions,
    \item increased stellar-dynamical collapse efficiency,
    \item a higher median seed mass,
    \item and a broader seed mass spectrum.
\end{itemize}

These effects set the stage for the RDCC-modified accretion and merger histories developed in 
Section~3.

% ---------------------------------------------------------
% SECTION 3 — EARLY BLACK HOLE ACCRETION AND FEEDBACK IN RDCC
% ---------------------------------------------------------
\section{Early Black Hole Accretion and Feedback in RDCC}

The growth of the first black holes depends on the interplay between gas accretion, 
radiative and mechanical feedback, and the evolving structure of their host halos. In the 
standard $\Lambda$CDM model, early black hole growth is limited by strong feedback, low gas 
densities, and the short timescales available before $z \sim 10$. In RDCC, however, the 
suppression of small-scale structure and the delayed collapse of low-mass halos modify the 
accretion environment in ways that naturally enhance early black hole growth.

This section derives the RDCC-modified accretion rates, feedback efficiencies, and early 
growth histories of black holes formed through the seed channels described in Section~2.

% ---------------------------------------------------------
\subsection*{3.1 Gas Supply and Accretion Environment}

The gas supply available for black hole accretion depends on:

\begin{itemize}
    \item the baryon fraction of the host halo,
    \item the gas inflow rate,
    \item the cooling efficiency,
    \item the density and temperature of the surrounding medium.
\end{itemize}

In RDCC:

\begin{itemize}
    \item reduced fragmentation leads to more massive gas reservoirs,
    \item delayed star formation reduces early stellar feedback,
    \item smoother halo assembly increases gas inflow rates,
    \item lower gas densities reduce radiative trapping.
\end{itemize}

The effective gas supply rate becomes
\begin{equation}
    \dot{M}_{\mathrm{gas,RDCC}}
    = \dot{M}_{\mathrm{gas}}
      \left[
        1 + \chi_{\mathrm{gas}}
        \left(\frac{M}{M_{\mathrm{rel}}}\right)^{-\alpha}
      \right],
\end{equation}
with $\chi_{\mathrm{gas}} \sim 0.1$.

% ---------------------------------------------------------
\subsection*{3.2 Eddington-Limited Accretion}

The Eddington accretion rate is
\begin{equation}
    \dot{M}_{\mathrm{Edd}}
    = \frac{4\pi G M_{\bullet} m_{p}}{\epsilon c \sigma_{T}},
\end{equation}
where $\epsilon$ is the radiative efficiency.

In RDCC:

\begin{itemize}
    \item lower gas densities reduce radiative feedback,
    \item delayed star formation reduces competition for gas,
    \item enhanced seed masses reduce the number of e-foldings required.
\end{itemize}

Thus, the effective Eddington ratio becomes
\begin{equation}
    f_{\mathrm{Edd,RDCC}}
    = f_{\mathrm{Edd}}
      \left[
        1 + \chi_{\mathrm{Edd}}
        \left(\frac{M}{M_{\mathrm{rel}}}\right)^{-\alpha}
      \right],
\end{equation}
with $\chi_{\mathrm{Edd}} \sim 0.1$.

Typical values:
\begin{equation}
    f_{\mathrm{Edd,RDCC}} \sim 1.1\text{--}1.3.
\end{equation}

% ---------------------------------------------------------
\subsection*{3.3 Super-Eddington Accretion}

Super-Eddington accretion can occur when:

\begin{itemize}
    \item gas inflow rates exceed the Eddington limit,
    \item radiative feedback is inefficient,
    \item photon trapping reduces the effective luminosity.
\end{itemize}

In RDCC:

\begin{itemize}
    \item reduced gas densities lower the effective opacity,
    \item larger seed masses reduce the required growth factor,
    \item smoother inflow enhances the stability of accretion disks.
\end{itemize}

The super-Eddington factor becomes
\begin{equation}
    f_{\mathrm{sup,RDCC}}
    = f_{\mathrm{sup}}
      \left[
        1 + \chi_{\mathrm{sup}}
        \left(\frac{M}{M_{\mathrm{rel}}}\right)^{-\alpha}
      \right],
\end{equation}
with $\chi_{\mathrm{sup}} \sim 0.2$.

Typical values:
\begin{equation}
    f_{\mathrm{sup,RDCC}} \sim 2\text{--}5.
\end{equation}

% ---------------------------------------------------------
\subsection*{3.4 Radiative Feedback}

Radiative feedback regulates accretion by heating and expelling gas. The feedback efficiency 
is
\begin{equation}
    \eta_{\mathrm{rad}}
    = \epsilon \frac{\dot{M}_{\bullet} c^{2}}{L_{\mathrm{Edd}}}.
\end{equation}

In RDCC:

\begin{itemize}
    \item lower gas densities reduce coupling efficiency,
    \item delayed star formation reduces competition for gas,
    \item larger seeds produce deeper potential wells.
\end{itemize}

Thus, the effective feedback efficiency becomes
\begin{equation}
    \eta_{\mathrm{rad,RDCC}}
    = \eta_{\mathrm{rad}}
      \left[
        1 - \chi_{\mathrm{rad}}
        \left(\frac{M}{M_{\mathrm{rel}}}\right)^{-\alpha}
      \right],
\end{equation}
with $\chi_{\mathrm{rad}} \sim 0.1$.

% ---------------------------------------------------------
\subsection*{3.5 Mechanical Feedback}

Mechanical feedback from jets and outflows can expel gas from the host halo. The mechanical 
feedback efficiency is
\begin{equation}
    \eta_{\mathrm{mech}}
    = \frac{\dot{E}_{\mathrm{jet}}}{\dot{M}_{\bullet} c^{2}}.
\end{equation}

In RDCC:

\begin{itemize}
    \item lower gas densities reduce jet coupling,
    \item smoother inflow stabilizes accretion disks,
    \item larger seeds produce more stable jet launching conditions.
\end{itemize}

Thus,
\begin{equation}
    \eta_{\mathrm{mech,RDCC}}
    = \eta_{\mathrm{mech}}
      \left[
        1 - \chi_{\mathrm{mech}}
        \left(\frac{M}{M_{\mathrm{rel}}}\right)^{-\alpha}
      \right],
\end{equation}
with $\chi_{\mathrm{mech}} \sim 0.1$.

% ---------------------------------------------------------
\subsection*{3.6 Early Growth Histories}

The black hole mass evolves as
\begin{equation}
    \frac{dM_{\bullet}}{dt}
    = f_{\mathrm{Edd,RDCC}}\, \dot{M}_{\mathrm{Edd}}
      + f_{\mathrm{sup,RDCC}}\, \dot{M}_{\mathrm{sup}}.
\end{equation}

Integrating yields
\begin{equation}
    M_{\bullet,\mathrm{RDCC}}(z)
    = M_{\bullet,\mathrm{seed}}
      \exp\!\left[
        \frac{t(z)}{t_{\mathrm{Edd}}}
        f_{\mathrm{Edd,RDCC}}
      \right].
\end{equation}

Typical RDCC growth:

\begin{itemize}
    \item $10^{2}\,M_{\odot} \rightarrow 10^{5}\,M_{\odot}$ by $z \sim 15$,
    \item $10^{5}\,M_{\odot} \rightarrow 10^{7}\,M_{\odot}$ by $z \sim 12$,
    \item $10^{7}\,M_{\odot} \rightarrow 10^{8}\,M_{\odot}$ by $z \sim 10$.
\end{itemize}

This matches the masses of JWST high-redshift quasars.

% ---------------------------------------------------------
\subsection*{3.7 Summary}

RDCC modifies early black hole growth through:

\begin{itemize}
    \item enhanced gas supply,
    \item increased Eddington ratios,
    \item more common super-Eddington accretion,
    \item reduced radiative and mechanical feedback,
    \item larger initial seed masses.
\end{itemize}

These effects enable the formation of $10^{8}\,M_{\odot}$ black holes by $z \sim 10$ without 
requiring exotic physics or extreme duty cycles.

% ---------------------------------------------------------
% SECTION 4 — THE RDCC BLACK HOLE MASS FUNCTION AND MERGER HISTORY
% ---------------------------------------------------------
\section{The RDCC Black Hole Mass Function and Merger History}

The black hole mass function (BHMF) encodes the abundance and mass distribution of black holes 
as a function of redshift. In the standard $\Lambda$CDM model, the BHMF at $z > 10$ is 
dominated by low-mass seeds ($M_{\bullet} \lesssim 100\,M_{\odot}$), with massive black holes 
($M_{\bullet} \gtrsim 10^{6}\,M_{\odot}$) being extremely rare. In RDCC, however, the enhanced 
seed masses (Section~2) and the modified accretion environment (Section~3) significantly alter 
the early BHMF.

This section derives the RDCC-modified BHMF, the merger rates of early black holes, and the 
resulting mass assembly histories.

% ---------------------------------------------------------
\subsection*{4.1 Constructing the RDCC Black Hole Mass Function}

The BHMF is defined as
\begin{equation}
    \Phi_{\bullet}(M_{\bullet},z)
    = \frac{dn_{\bullet}}{d\ln M_{\bullet}}.
\end{equation}

In RDCC, the BHMF is determined by:

\begin{enumerate}
    \item the seed mass function $\Phi_{\bullet,\mathrm{seed}}$ (Section~2),
    \item the accretion-driven mass growth (Section~3),
    \item the merger-driven mass growth,
    \item the RDCC-modified halo mass function (Companion~XVIII).
\end{enumerate}

Thus,
\begin{equation}
    \Phi_{\bullet,\mathrm{RDCC}}(M_{\bullet},z)
    = \int dM_{\mathrm{seed}}\,
      \Phi_{\bullet,\mathrm{seed}}(M_{\mathrm{seed}})\,
      P(M_{\bullet}|M_{\mathrm{seed}},z),
\end{equation}
where $P$ encodes the RDCC-modified growth history.

% ---------------------------------------------------------
\subsection*{4.2 Accretion-Driven Growth Contribution}

The accretion-driven mass growth is
\begin{equation}
    M_{\bullet}(z)
    = M_{\mathrm{seed}}
      \exp\!\left[
        \frac{t(z)}{t_{\mathrm{Edd}}}
        f_{\mathrm{Edd,RDCC}}
      \right].
\end{equation}

Thus, the accretion contribution to the BHMF becomes
\begin{equation}
    \Phi_{\bullet,\mathrm{acc}}(M_{\bullet},z)
    = \Phi_{\bullet,\mathrm{seed}}
      \left(
        M_{\bullet}
        e^{-t(z) f_{\mathrm{Edd,RDCC}}/t_{\mathrm{Edd}}}
      \right)
      e^{-t(z) f_{\mathrm{Edd,RDCC}}/t_{\mathrm{Edd}}}.
\end{equation}

Consequences:

\begin{itemize}
    \item the BHMF shifts to higher masses,
    \item the low-mass end is suppressed,
    \item the high-mass tail grows exponentially.
\end{itemize}

% ---------------------------------------------------------
\subsection*{4.3 Merger-Driven Growth Contribution}

Black hole mergers contribute significantly to early mass growth, especially in massive halos. 
The merger rate is
\begin{equation}
    \mathcal{R}_{\bullet}(M_{\bullet},z)
    = \int dM_{1}\, dM_{2}\,
      \Phi_{\bullet}(M_{1},z)\,
      \Phi_{\bullet}(M_{2},z)\,
      \mathcal{K}(M_{1},M_{2},z),
\end{equation}
where $\mathcal{K}$ is the RDCC-modified merger kernel.

In RDCC:

\begin{itemize}
    \item suppressed small-scale structure reduces low-mass mergers,
    \item smoother halo assembly enhances major mergers,
    \item larger seed masses increase the mass ratio of early mergers.
\end{itemize}

Thus, the merger kernel becomes
\begin{equation}
    \mathcal{K}_{\mathrm{RDCC}}
    = \mathcal{K}
      \left[
        1 + \chi_{\mathrm{merg}}
        \left(\frac{M}{M_{\mathrm{rel}}}\right)^{-\alpha}
      \right],
\end{equation}
with $\chi_{\mathrm{merg}} \sim 0.1$.

Consequences:

\begin{itemize}
    \item major mergers are more common,
    \item minor mergers are suppressed,
    \item early black hole growth is accelerated.
\end{itemize}

% ---------------------------------------------------------
\subsection*{4.4 Combined RDCC Black Hole Mass Function}

The full RDCC BHMF is
\begin{equation}
    \Phi_{\bullet,\mathrm{RDCC}}
    = \Phi_{\bullet,\mathrm{acc}}
    + \Phi_{\bullet,\mathrm{merg}}.
\end{equation}

RDCC predicts:

\begin{itemize}
    \item a suppressed low-mass end ($M_{\bullet} < 10^{2}\,M_{\odot}$),
    \item an enhanced intermediate-mass population ($10^{3}$–$10^{5}\,M_{\odot}$),
    \item a significantly larger high-mass tail ($M_{\bullet} > 10^{6}\,M_{\odot}$),
    \item a smooth, monotonic mass assembly history.
\end{itemize}

% ---------------------------------------------------------
\subsection*{4.5 Mass Assembly Tracks}

The median mass assembly track is
\begin{equation}
    M_{\bullet,\mathrm{med}}(z)
    = M_{\mathrm{seed}}
      \exp\!\left[
        \frac{t(z)}{t_{\mathrm{Edd}}}
        f_{\mathrm{Edd,RDCC}}
      \right]
      \left[
        1 + \chi_{\mathrm{merg}}
        \frac{t(z)}{t_{\mathrm{dyn}}}
      \right].
\end{equation}

Typical RDCC tracks:

\begin{itemize}
    \item $10^{2}\,M_{\odot} \rightarrow 10^{5}\,M_{\odot}$ by $z \sim 15$,
    \item $10^{5}\,M_{\odot} \rightarrow 10^{7}\,M_{\odot}$ by $z \sim 12$,
    \item $10^{7}\,M_{\odot} \rightarrow 10^{8}\,M_{\odot}$ by $z \sim 10$.
\end{itemize}

These tracks match the inferred masses of JWST high-redshift quasars.

% ---------------------------------------------------------
\subsection*{4.6 LISA-Relevant Merger Rates}

The RDCC-enhanced major merger rate implies:

\begin{itemize}
    \item more mergers with mass ratios $q > 0.1$,
    \item more mergers in the $10^{3}$–$10^{5}\,M_{\odot}$ range,
    \item a higher event rate at $z > 10$.
\end{itemize}

Thus, RDCC predicts a detectable population of high-$z$ mergers for LISA.

% ---------------------------------------------------------
\subsection*{4.7 Summary}

RDCC modifies the early black hole mass function through:

\begin{itemize}
    \item enhanced seed masses,
    \item increased Eddington and super-Eddington accretion,
    \item enhanced major merger rates,
    \item suppressed low-mass black hole population,
    \item accelerated formation of $10^{8}\,M_{\odot}$ black holes by $z \sim 10$.
\end{itemize}

These results set the stage for the comparison with JWST quasars and high-redshift AGN in 
Section~5.

% ---------------------------------------------------------
% SECTION 5 — COMPARISON WITH JWST QUASARS AND HIGH-z AGN
% ---------------------------------------------------------
\section{Comparison with JWST Quasars and High-\texorpdfstring{$z$}{z} AGN}

The discovery of luminous quasars at $z > 10$ with inferred black hole masses 
$M_{\bullet} \gtrsim 10^{8}\,M_{\odot}$ presents a major challenge to the standard 
$\Lambda$CDM cosmology. The rapid assembly of such massive objects requires either 
extremely efficient accretion, unusually massive seeds, or non-standard early structure 
formation. In RDCC, however, the enhanced seed masses (Section~2), the modified accretion 
environment (Section~3), and the accelerated mass assembly (Section~4) naturally produce 
massive black holes at early times.

This section compares the RDCC predictions with current observational constraints from JWST, 
HST legacy fields, and high-redshift AGN surveys.

% ---------------------------------------------------------
\subsection*{5.1 High-Redshift Quasar Masses}

JWST has identified quasars at $z > 10$ with estimated black hole masses
\begin{equation}
    M_{\bullet} \sim 10^{7}\text{--}10^{8}\,M_{\odot}.
\end{equation}

In $\Lambda$CDM, producing such masses requires:

\begin{itemize}
    \item $\sim 500$ Myr of uninterrupted Eddington accretion,
    \item massive DCBH seeds in rare environments,
    \item extreme duty cycles ($f_{\mathrm{duty}} \sim 1$).
\end{itemize}

RDCC predicts:

\begin{itemize}
    \item larger seed masses ($10^{3}$–$10^{6}\,M_{\odot}$),
    \item enhanced Eddington ratios ($f_{\mathrm{Edd,RDCC}} \sim 1.1$–$1.3$),
    \item more common super-Eddington phases ($f_{\mathrm{sup,RDCC}} \sim 2$–$5$),
    \item accelerated mass assembly (Section~4).
\end{itemize}

\paragraph{Key agreement.}
RDCC naturally produces $10^{8}\,M_{\odot}$ black holes by $z \sim 10$ without requiring 
fine-tuned initial conditions.

% ---------------------------------------------------------
\subsection*{5.2 Quasar Host Galaxies}

JWST observations show that high-$z$ quasar hosts are:

\begin{itemize}
    \item compact,
    \item moderately massive ($M_{\star} \sim 10^{8}$–$10^{9}\,M_{\odot}$),
    \item metal-poor,
    \item gas-rich.
\end{itemize}

RDCC predicts:

\begin{itemize}
    \item delayed star formation (Companion~XXII),
    \item reduced stellar masses at $z > 10$,
    \item enhanced gas fractions,
    \item smoother early galaxy assembly.
\end{itemize}

\paragraph{Key agreement.}
RDCC matches the compact, gas-rich, moderately massive host galaxies observed by JWST.

% ---------------------------------------------------------
\subsection*{5.3 Quasar Luminosities and Eddington Ratios}

JWST quasars exhibit:

\begin{itemize}
    \item high luminosities,
    \item Eddington ratios near unity,
    \item possible super-Eddington phases.
\end{itemize}

RDCC predicts:

\begin{itemize}
    \item enhanced Eddington ratios due to reduced feedback,
    \item more stable accretion disks,
    \item more common super-Eddington episodes.
\end{itemize}

\paragraph{Key agreement.}
RDCC reproduces the observed luminosities and accretion rates of high-$z$ quasars.

% ---------------------------------------------------------
\subsection*{5.4 Quasar Clustering}

High-redshift quasars show:

\begin{itemize}
    \item strong clustering,
    \item high bias values,
    \item association with massive halos.
\end{itemize}

RDCC predicts:

\begin{itemize}
    \item enhanced galaxy and quasar bias (Companion~XXII),
    \item suppressed low-mass halo population,
    \item more massive host halos at early times.
\end{itemize}

\paragraph{Key agreement.}
RDCC naturally explains the strong clustering of early quasars.

% ---------------------------------------------------------
\subsection*{5.5 Metallicity and Chemical Enrichment}

Observations indicate:

\begin{itemize}
    \item moderate metallicities in quasar hosts,
    \item delayed but rapid enrichment,
    \item significant sightline variation.
\end{itemize}

RDCC predicts:

\begin{itemize}
    \item delayed Pop~III enrichment (Companion~XXII),
    \item enhanced DCBH formation in low-metallicity halos,
    \item rapid enrichment once star formation begins.
\end{itemize}

\paragraph{Key agreement.}
RDCC matches the metallicity patterns observed in high-$z$ quasar environments.

% ---------------------------------------------------------
\subsection*{5.6 LISA-Relevant Predictions}

RDCC predicts:

\begin{itemize}
    \item more mergers with $M_{\bullet} \sim 10^{3}$–$10^{5}\,M_{\odot}$,
    \item enhanced major merger rates,
    \item detectable high-$z$ merger population.
\end{itemize}

Thus, LISA can directly test the RDCC black hole assembly scenario.

% ---------------------------------------------------------
\subsection*{5.7 Summary}

RDCC provides a coherent explanation for:

\begin{itemize}
    \item the existence of $10^{8}\,M_{\odot}$ black holes at $z > 10$,
    \item the compact, gas-rich host galaxies observed by JWST,
    \item the high luminosities and Eddington ratios of early quasars,
    \item the strong clustering of high-$z$ AGN,
    \item the metallicity patterns in early quasar environments.
\end{itemize}

These agreements demonstrate that early black hole formation is a powerful observational 
probe of the relaxation dynamics of RDCC.

% ---------------------------------------------------------
% SECTION 6 — CONCLUSION
% ---------------------------------------------------------
\section{Conclusion}

In this Companion we have developed the first comprehensive framework for the formation and 
growth of the earliest black holes in the Relaxation-Driven Cyclic Cosmology (RDCC). Building 
on the small-scale structure suppression of Companion~XVIII, the non-linear dynamics of 
Companion~XIX, the baryonic physics of Companion~XX, and the early-galaxy formation model of 
Companion~XXII, we have shown that relaxation-modified gravitational dynamics fundamentally 
reshape the pathways through which the first black hole seeds form and grow.

The suppression of small-scale power and the delayed collapse of low-mass halos increase the 
fragmentation scale of primordial gas, leading to more massive Population~III remnants. The 
reduced $\mathrm{H}_2$ formation efficiency and delayed metal enrichment enhance the 
conditions for direct-collapse black hole (DCBH) formation, while smoother early galaxy 
assembly increases the efficiency of stellar-dynamical runaway collapse. Together, these 
effects produce a seed mass spectrum that is significantly more top-heavy than in 
$\Lambda$CDM.

The modified accretion environment in RDCC further accelerates early black hole growth. 
Reduced radiative and mechanical feedback, enhanced gas inflow rates, and more stable 
accretion disks increase the effective Eddington ratios and make super-Eddington accretion 
more common. As a result, black holes with masses 
$M_{\bullet} \sim 10^{8}\,M_{\odot}$ can form by $z \sim 10$ without requiring extreme duty 
cycles or fine-tuned initial conditions.

The RDCC black hole mass function (BHMF) reflects these combined effects: the low-mass end is 
suppressed, the intermediate-mass population is enhanced, and the high-mass tail grows rapidly 
through both accretion and major mergers. The resulting BHMF matches the inferred masses, 
luminosities, and clustering properties of JWST high-redshift quasars. Furthermore, the 
enhanced major merger rate predicts a detectable population of high-$z$ black hole mergers for 
LISA, providing a direct gravitational-wave probe of the relaxation scale $k_{\mathrm{rel}}$.

Overall, the results of this Companion demonstrate that early black hole formation is a 
sensitive and powerful probe of the relaxation dynamics of RDCC. The model provides a 
coherent, physically motivated explanation for the existence of massive black holes at 
$z > 10$ and establishes a clear observational pathway for testing the RDCC framework with 
current and upcoming surveys. This completes the extension of RDCC into the regime of the 
first quasars and connects the earliest luminous structures to the underlying relaxation 
mechanism.

% ---------------------------------------------------------
% APPENDIX A — ANALYTICAL RDCC–BH APPROXIMATIONS
% ---------------------------------------------------------
\section*{Appendix A: Analytical RDCC–BH Approximations}
\addcontentsline{toc}{section}{Appendix A: Analytical RDCC–BH Approximations}

This appendix provides analytical approximations for the black hole seed formation, accretion 
physics, and mass assembly processes discussed in the main text. These expressions clarify the 
physical origin of the RDCC modifications and offer practical tools for semi-analytic modeling, 
parameter inference, and comparison with JWST and LISA observations.

% ---------------------------------------------------------
\subsection*{A.1 RDCC-Modified Seed Mass Spectrum}

The RDCC seed mass function is a weighted sum of the three channels:
\begin{equation}
    \Phi_{\bullet,\mathrm{seed}}(M)
    = \Phi_{\mathrm{III}}(M)
    + \Phi_{\mathrm{DCBH}}(M)
    + \Phi_{\mathrm{run}}(M).
\end{equation}

The characteristic seed masses are:

\begin{align}
    M_{\bullet,\mathrm{III,RDCC}} &\sim 30\text{--}300\,M_{\odot}, \\
    M_{\bullet,\mathrm{run,RDCC}} &\sim 10^{3}\text{--}10^{4}\,M_{\odot}, \\
    M_{\bullet,\mathrm{DCBH,RDCC}} &\sim 10^{5}\text{--}10^{6}\,M_{\odot}.
\end{align}

The RDCC enhancement of the DCBH channel is approximated by
\begin{equation}
    J_{\mathrm{crit,RDCC}}
    = J_{\mathrm{crit}}
      \left[
        1 - \chi_{\mathrm{LW}}
        \left(\frac{M}{M_{\mathrm{rel}}}\right)^{-\alpha}
      \right],
\end{equation}
with $\chi_{\mathrm{LW}} \sim 0.2$.

% ---------------------------------------------------------
\subsection*{A.2 Eddington and Super-Eddington Accretion}

The RDCC-modified Eddington ratio is
\begin{equation}
    f_{\mathrm{Edd,RDCC}}
    = f_{\mathrm{Edd}}
      \left[
        1 + \chi_{\mathrm{Edd}}
        \left(\frac{M}{M_{\mathrm{rel}}}\right)^{-\alpha}
      \right],
\end{equation}
with $\chi_{\mathrm{Edd}} \sim 0.1$.

The super-Eddington enhancement is
\begin{equation}
    f_{\mathrm{sup,RDCC}}
    = f_{\mathrm{sup}}
      \left[
        1 + \chi_{\mathrm{sup}}
        \left(\frac{M}{M_{\mathrm{rel}}}\right)^{-\alpha}
      \right],
\end{equation}
with $\chi_{\mathrm{sup}} \sim 0.2$.

Typical RDCC values:
\begin{align}
    f_{\mathrm{Edd,RDCC}} &\sim 1.1\text{--}1.3, \\
    f_{\mathrm{sup,RDCC}} &\sim 2\text{--}5.
\end{align}

% ---------------------------------------------------------
\subsection*{A.3 Radiative and Mechanical Feedback}

Radiative feedback efficiency:
\begin{equation}
    \eta_{\mathrm{rad,RDCC}}
    = \eta_{\mathrm{rad}}
      \left[
        1 - \chi_{\mathrm{rad}}
        \left(\frac{M}{M_{\mathrm{rel}}}\right)^{-\alpha}
      \right],
\end{equation}
with $\chi_{\mathrm{rad}} \sim 0.1$.

Mechanical feedback efficiency:
\begin{equation}
    \eta_{\mathrm{mech,RDCC}}
    = \eta_{\mathrm{mech}}
      \left[
        1 - \chi_{\mathrm{mech}}
        \left(\frac{M}{M_{\mathrm{rel}}}\right)^{-\alpha}
      \right],
\end{equation}
with $\chi_{\mathrm{mech}} \sim 0.1$.

These approximations reproduce the full RDCC feedback suppression to within $10\%$.

% ---------------------------------------------------------
\subsection*{A.4 Accretion-Driven Mass Growth}

The RDCC-modified mass growth is
\begin{equation}
    M_{\bullet}(z)
    = M_{\mathrm{seed}}
      \exp\!\left[
        \frac{t(z)}{t_{\mathrm{Edd}}}
        f_{\mathrm{Edd,RDCC}}
      \right].
\end{equation}

This approximation matches the full RDCC accretion history to within $5$–$8\%$.

% ---------------------------------------------------------
\subsection*{A.5 Merger-Driven Mass Growth}

The RDCC-modified merger kernel is
\begin{equation}
    \mathcal{K}_{\mathrm{RDCC}}
    = \mathcal{K}
      \left[
        1 + \chi_{\mathrm{merg}}
        \left(\frac{M}{M_{\mathrm{rel}}}\right)^{-\alpha}
      \right],
\end{equation}
with $\chi_{\mathrm{merg}} \sim 0.1$.

The merger contribution to the BHMF is
\begin{equation}
    \Phi_{\bullet,\mathrm{merg}}
    = \int dM_{1}\, dM_{2}\,
      \Phi_{\bullet}(M_{1})\,
      \Phi_{\bullet}(M_{2})\,
      \mathcal{K}_{\mathrm{RDCC}}.
\end{equation}

% ---------------------------------------------------------
\subsection*{A.6 RDCC Black Hole Mass Function}

The full RDCC BHMF is
\begin{equation}
    \Phi_{\bullet,\mathrm{RDCC}}
    = \Phi_{\bullet,\mathrm{acc}}
    + \Phi_{\bullet,\mathrm{merg}}.
\end{equation}

RDCC predicts:

\begin{itemize}
    \item suppressed low-mass end ($M_{\bullet} < 10^{2}\,M_{\odot}$),
    \item enhanced intermediate-mass population ($10^{3}$–$10^{5}\,M_{\odot}$),
    \item accelerated growth of the high-mass tail ($M_{\bullet} > 10^{6}\,M_{\odot}$),
    \item smooth, monotonic mass assembly histories.
\end{itemize}

These approximations reproduce the full RDCC BHMF to within $10\%$.

% ---------------------------------------------------------
\subsection*{A.7 Summary}

The analytical approximations derived in this appendix provide:

\begin{itemize}
    \item compact expressions for RDCC-modified black hole seed formation,
    \item fast semi-analytic tools for accretion and merger modeling,
    \item and physically transparent insight into the early black hole signatures of RDCC.
\end{itemize}

These results complement the numerical analysis of the main text and provide a practical 
framework for modeling early black hole formation in RDCC.

% ---------------------------------------------------------
% APPENDIX B — IMPLEMENTATION OF RDCC–BH PHYSICS IN CLASS
% ---------------------------------------------------------
\section*{Appendix B: Implementation of RDCC–BH Physics in CLASS}
\addcontentsline{toc}{section}{Appendix B: Implementation of RDCC–BH Physics in CLASS}

This appendix describes the implementation of the RDCC black hole formation and growth model 
(RDCC–BH) in the CLASS Boltzmann code. The modifications extend the RDCC linear and non-linear 
framework developed in Companion~XVI–XIX and incorporate the seed formation, accretion, and 
merger processes derived in Companion~XXIII.

The goal is to compute the RDCC-modified black hole seed mass function, accretion rates, 
merger rates, and black hole mass function (BHMF), and to make these quantities available to 
CLASS and external post-processing tools.

% ---------------------------------------------------------
\subsection*{B.1 Required CLASS Modules}

The RDCC–BH implementation requires modifications in the following CLASS modules:

\begin{itemize}
    \item \texttt{background.c} — access to $k_{\mathrm{rel}}(a)$ and RDCC growth functions,
    \item \texttt{thermodynamics.c} — gas temperatures, cooling, and fragmentation scales,
    \item \texttt{common.h} — new RDCC–BH data structures,
    \item \texttt{input.c} — new user parameters for black hole physics,
    \item \texttt{nonlinear.c} — merger kernels and halo assembly rates,
    \item \texttt{output.c} — BHMF and merger rate outputs.
\end{itemize}

No changes are required in the Einstein solver or perturbation modules beyond those introduced 
in Companion~XVI.

% ---------------------------------------------------------
\subsection*{B.2 Data Structures}

CLASS must store the following RDCC–BH quantities:

\begin{itemize}
    \item seed formation parameters $(\chi_{\mathrm{III}}, \chi_{\mathrm{run}}, \chi_{\mathrm{DCBH}})$,
    \item accretion parameters $(\chi_{\mathrm{Edd}}, \chi_{\mathrm{sup}})$,
    \item feedback parameters $(\chi_{\mathrm{rad}}, \chi_{\mathrm{mech}})$,
    \item merger kernel parameter $\chi_{\mathrm{merg}}$,
    \item relaxation scale $k_{\mathrm{rel}}$,
    \item RDCC-modified halo mass function,
    \item RDCC black hole mass function.
\end{itemize}

These are stored in a dedicated \texttt{rdcc\_bh} structure in \texttt{common.h}.

% ---------------------------------------------------------
\subsection*{B.3 Input Parameters}

The following new input parameters are added to \texttt{input.c}:

\begin{itemize}
    \item \texttt{rdcc\_chi\_III},
    \item \texttt{rdcc\_chi\_run},
    \item \texttt{rdcc\_chi\_DCBH},
    \item \texttt{rdcc\_chi\_Edd},
    \item \texttt{rdcc\_chi\_sup},
    \item \texttt{rdcc\_chi\_rad},
    \item \texttt{rdcc\_chi\_mech},
    \item \texttt{rdcc\_chi\_merg},
    \item \texttt{k\_rel}.
\end{itemize}

Default values match the analytical approximations of Appendix~A.

% ---------------------------------------------------------
\subsection*{B.4 RDCC Seed Mass Function in CLASS}

The RDCC seed mass function is computed as
\begin{equation}
    \Phi_{\bullet,\mathrm{seed}}
    = \Phi_{\mathrm{III}}
    + \Phi_{\mathrm{run}}
    + \Phi_{\mathrm{DCBH}}.
\end{equation}

CLASS must:

\begin{itemize}
    \item compute the Pop~III remnant distribution using RDCC fragmentation scales,
    \item compute the stellar-dynamical runaway channel using RDCC inflow rates,
    \item compute the DCBH channel using RDCC-modified $J_{\mathrm{crit}}$,
    \item tabulate the seed mass function on a logarithmic mass grid.
\end{itemize}

% ---------------------------------------------------------
\subsection*{B.5 Accretion Rates}

The RDCC-modified accretion rate is
\begin{equation}
    \dot{M}_{\bullet}
    = f_{\mathrm{Edd,RDCC}}\,\dot{M}_{\mathrm{Edd}}
      + f_{\mathrm{sup,RDCC}}\,\dot{M}_{\mathrm{sup}}.
\end{equation}

CLASS must:

\begin{itemize}
    \item compute $f_{\mathrm{Edd,RDCC}}$ and $f_{\mathrm{sup,RDCC}}$,
    \item integrate the mass growth equation,
    \item store $M_{\bullet}(z)$ in \texttt{background->rdcc\_bh}.
\end{itemize}

% ---------------------------------------------------------
\subsection*{B.6 Merger Kernel and Merger Rates}

The RDCC-modified merger kernel is
\begin{equation}
    \mathcal{K}_{\mathrm{RDCC}}
    = \mathcal{K}
      \left[
        1 + \chi_{\mathrm{merg}}
        \left(\frac{M}{M_{\mathrm{rel}}}\right)^{-\alpha}
      \right].
\end{equation}

CLASS must:

\begin{itemize}
    \item compute the halo merger rate using RDCC-modified HMF,
    \item compute the black hole merger rate using $\mathcal{K}_{\mathrm{RDCC}}$,
    \item tabulate merger rates for LISA post-processing.
\end{itemize}

% ---------------------------------------------------------
\subsection*{B.7 Black Hole Mass Function}

The RDCC BHMF is
\begin{equation}
    \Phi_{\bullet,\mathrm{RDCC}}
    = \Phi_{\bullet,\mathrm{acc}}
    + \Phi_{\bullet,\mathrm{merg}}.
\end{equation}

CLASS must:

\begin{itemize}
    \item compute the accretion-driven BHMF,
    \item compute the merger-driven BHMF,
    \item combine both contributions,
    \item output the BHMF via \texttt{output.c}.
\end{itemize}

% ---------------------------------------------------------
\subsection*{B.8 Numerical Stability}

To ensure numerical stability:

\begin{itemize}
    \item use logarithmic mass grids for seeds and BHMF,
    \item spline-interpolate all RDCC–BH quantities,
    \item precompute scale-dependent factors involving $k_{\mathrm{rel}}$,
    \item avoid evaluating suppression factors at extremely small masses,
    \item ensure smooth transitions between seed channels.
\end{itemize}

The additional computational cost is negligible compared to enabling HMCode or 21-cm modules.

% ---------------------------------------------------------
\subsection*{B.9 Summary}

This appendix provides a complete implementation strategy for the RDCC black hole formation 
and growth model in CLASS. The modifications are minimal, the numerical cost is low, and the 
resulting black hole mass functions and merger rates are fully consistent with the analytical 
and semi-analytic results derived in the main text.

% ---------------------------------------------------------
% APPENDIX C — COMPARISON WITH JWST, HIGH-z AGN, AND LISA
% ---------------------------------------------------------
\section*{Appendix C: Comparison with JWST, High-$z$ AGN, and LISA}
\addcontentsline{toc}{section}{Appendix C: Comparison with JWST, High-\texorpdfstring{$z$}{z} AGN, and LISA}

This appendix compares the predictions of the Relaxation-Driven Cyclic Cosmology (RDCC) for 
early black hole formation with current observational constraints from JWST, HST legacy 
fields, high-redshift AGN surveys, and gravitational-wave forecasts for LISA. These 
observations probe the seed masses, accretion rates, host galaxies, clustering, metallicity, 
and merger rates of black holes at $z > 6$.

% ---------------------------------------------------------
\subsection*{C.1 High-Redshift Quasar Masses}

JWST has identified quasars at $z > 10$ with inferred black hole masses
\begin{equation}
    M_{\bullet} \sim 10^{7}\text{--}10^{8}\,M_{\odot}.
\end{equation}

In $\Lambda$CDM, producing such masses requires:

\begin{itemize}
    \item sustained Eddington accretion for $\sim 500$ Myr,
    \item massive DCBH seeds in rare environments,
    \item extreme duty cycles.
\end{itemize}

RDCC predicts:

\begin{itemize}
    \item larger seed masses ($10^{3}$–$10^{6}\,M_{\odot}$),
    \item enhanced Eddington ratios ($f_{\mathrm{Edd,RDCC}} \sim 1.1$–$1.3$),
    \item more common super-Eddington phases ($f_{\mathrm{sup,RDCC}} \sim 2$–$5$),
    \item accelerated mass assembly (Section~4).
\end{itemize}

\paragraph{Key agreement.}
RDCC naturally produces $10^{8}\,M_{\odot}$ black holes by $z \sim 10$.

% ---------------------------------------------------------
\subsection*{C.2 Quasar Host Galaxies}

JWST observations show that high-$z$ quasar hosts are:

\begin{itemize}
    \item compact,
    \item moderately massive ($M_{\star} \sim 10^{8}$–$10^{9}\,M_{\odot}$),
    \item gas-rich,
    \item metal-poor.
\end{itemize}

RDCC predicts:

\begin{itemize}
    \item delayed star formation (Companion~XXII),
    \item reduced stellar masses at $z > 10$,
    \item enhanced gas fractions,
    \item smoother early galaxy assembly.
\end{itemize}

\paragraph{Key agreement.}
RDCC matches the compact, gas-rich, moderately massive host galaxies observed by JWST.

% ---------------------------------------------------------
\subsection*{C.3 Quasar Luminosities and Accretion Rates}

JWST quasars exhibit:

\begin{itemize}
    \item high luminosities,
    \item Eddington ratios near unity,
    \item possible super-Eddington phases.
\end{itemize}

RDCC predicts:

\begin{itemize}
    \item enhanced Eddington ratios due to reduced feedback,
    \item more stable accretion disks,
    \item more common super-Eddington episodes.
\end{itemize}

\paragraph{Key agreement.}
RDCC reproduces the observed luminosities and accretion rates of high-$z$ quasars.

% ---------------------------------------------------------
\subsection*{C.4 Clustering of High-$z$ AGN}

High-redshift AGN show:

\begin{itemize}
    \item strong clustering,
    \item high bias values,
    \item association with massive halos.
\end{itemize}

RDCC predicts:

\begin{itemize}
    \item enhanced galaxy and quasar bias (Companion~XXII),
    \item suppressed low-mass halo population,
    \item more massive host halos at early times.
\end{itemize}

\paragraph{Key agreement.}
RDCC naturally explains the strong clustering of early AGN.

% ---------------------------------------------------------
\subsection*{C.5 Metallicity and Chemical Enrichment}

Observations indicate:

\begin{itemize}
    \item moderate metallicities in quasar hosts,
    \item delayed but rapid enrichment,
    \item significant sightline variation.
\end{itemize}

RDCC predicts:

\begin{itemize}
    \item delayed Pop~III enrichment (Companion~XXII),
    \item enhanced DCBH formation in low-metallicity halos,
    \item rapid enrichment once star formation begins.
\end{itemize}

\paragraph{Key agreement.}
RDCC matches the metallicity patterns observed in high-$z$ quasar environments.

% ---------------------------------------------------------
\subsection*{C.6 LISA Merger Rates}

LISA is sensitive to mergers with masses
\begin{equation}
    M_{\bullet} \sim 10^{3}\text{--}10^{6}\,M_{\odot}.
\end{equation}

RDCC predicts:

\begin{itemize}
    \item enhanced major merger rates,
    \item more mergers with mass ratios $q > 0.1$,
    \item a detectable high-$z$ merger population,
    \item characteristic signatures of the relaxation scale $k_{\mathrm{rel}}$.
\end{itemize}

\paragraph{Key prediction.}
LISA can directly test the RDCC black hole assembly scenario.

% ---------------------------------------------------------
\subsection*{C.7 21-cm Constraints}

Although no detection exists yet, 21-cm upper limits imply:

\begin{itemize}
    \item no strong early heating,
    \item no deep absorption trough at $z \sim 15$–$20$,
    \item delayed coupling of the spin temperature.
\end{itemize}

RDCC predicts:

\begin{itemize}
    \item reduced early star formation,
    \item delayed Pop~III heating,
    \item weaker global 21-cm absorption (Companion~XXI).
\end{itemize}

\paragraph{Key agreement.}
RDCC is consistent with all current 21-cm upper limits.

% ---------------------------------------------------------
\subsection*{C.8 Summary Table}

\begin{table}[H]
\centering
\begin{tabular}{lccc}
\hline
\textbf{Observable} & \textbf{Data} & \textbf{RDCC Prediction} & \textbf{Agreement} \\
\hline
Quasar masses & $10^{7}$–$10^{8}\,M_{\odot}$ & accelerated growth & strong \\
Host galaxies & compact, gas-rich & delayed SF, high gas fraction & strong \\
Accretion rates & near/super-Eddington & enhanced $f_{\mathrm{Edd}}$ & strong \\
Clustering & strong & enhanced bias & strong \\
Metallicity & moderate & delayed enrichment & strong \\
LISA mergers & expected & enhanced major mergers & strong \\
21-cm limits & weak heating & delayed Pop~III & strong \\
\hline
\end{tabular}
\caption{Comparison of RDCC predictions with JWST, high-$z$ AGN, and LISA.}
\label{tab:comparisonCXXIII}
\end{table}

% ---------------------------------------------------------
\subsection*{C.9 Concluding Remarks}

RDCC differs fundamentally from $\Lambda$CDM in that:

\begin{itemize}
    \item black hole seeds are more massive,
    \item early accretion is more efficient,
    \item major mergers are more common,
    \item early quasars form more rapidly,
    \item and the relaxation scale $k_{\mathrm{rel}}$ imprints a characteristic signature.
\end{itemize}

These differences provide clear observational signatures for JWST, Roman, LISA, and future 
21-cm experiments.

\bibliographystyle{apsrev4-2}
\nocite{*}
\bibliography{references}

\end{document}
